% Self-contained LaTeX manuscript.
% Compile with: pdflatex affine_reverse_faber_krahn.tex
% Run pdflatex a second time to resolve all cross-references.
% No external bibliography file or figures are required.
\documentclass[11pt]{amsart}

\usepackage[T1]{fontenc}
\usepackage[utf8]{inputenc}
\usepackage{lmodern}
\usepackage[letterpaper,margin=1in]{geometry}
\usepackage{mathtools,amssymb}
\usepackage{microtype}
\usepackage{enumitem}
\usepackage{hyperref}
\hypersetup{
  pdftitle={An affine reverse Faber--Krahn inequality via information decay},
  pdfauthor={Matthew J. Colbrook and George Stepaniants},
  colorlinks=true,
  linkcolor=blue,
  citecolor=blue,
  urlcolor=blue
}
\urlstyle{same}
\allowdisplaybreaks[1]
\setlength{\emergencystretch}{2em}
\raggedbottom
\numberwithin{equation}{section}

\newtheorem{theorem}{Theorem}[section]
\newtheorem{proposition}[theorem]{Proposition}
\newtheorem{lemma}[theorem]{Lemma}
\newtheorem{corollary}[theorem]{Corollary}
\theoremstyle{definition}
\newtheorem{definition}[theorem]{Definition}
\theoremstyle{remark}
\newtheorem{remark}[theorem]{Remark}

\newcommand{\R}{\mathbb R}
\newcommand{\E}{\mathbb E}
\newcommand{\Prob}{\mathbb P}
\newcommand{\Id}{I_n}
\newcommand{\gam}{\gamma_n}
\newcommand{\CP}{C_{\mathrm P}}
\newcommand{\HS}{\mathrm{HS}}
\newcommand{\op}{\mathrm{op}}
\newcommand{\eps}{\varepsilon}
\newcommand{\norm}[1]{\lVert #1\rVert}
\DeclareMathOperator{\Cov}{Cov}
\DeclareMathOperator{\Var}{Var}
\DeclareMathOperator{\supp}{supp}
\DeclareMathOperator{\Div}{div}
\DeclareMathOperator{\tr}{tr}
\DeclareMathOperator{\SL}{SL}
\DeclareMathOperator{\Int}{int}

\title[Affine reverse Faber--Krahn]{An affine reverse Faber--Krahn inequality\\via information decay}
\author{Matthew J. Colbrook}
\address{Department of Applied Mathematics and Theoretical Physics,
University of Cambridge, Wilberforce Road, Cambridge CB3 0WA, United Kingdom}
\email{m.colbrook@damtp.cam.ac.uk}
\thanks{Matthew J. Colbrook is the corresponding author.}
\author{George Stepaniants}
\address{Department of Applied Mathematics and Theoretical Physics,
University of Cambridge, Wilberforce Road, Cambridge CB3 0WA, United Kingdom}
\curraddr{Courant Institute of Mathematical Sciences, New York University,
251 Mercer Street, New York, NY 10012, USA}
\email{gs852@cam.ac.uk}
\date{27 October 2026}
\subjclass[2020]{35P15, 52A23, 60F05, 60J65}
\keywords{Dirichlet eigenvalue, convex body, Gaussian position, Fisher information,
Stein kernel, Brownian bridge}

\begin{document}

\begin{abstract}
We prove that every convex body \(K\subset\R^n\) admits
\(S\in\SL_n(\R)\) such that
\(\lambda_1(SK)|K|^{2/n}\le Cn\), where \(C\) is universal.
For origin-symmetric bodies of volume one, the bound holds in maximal
Gaussian measure position. The main step is the spectral estimate
\[
\lambda_1(D)\le C(P)\bigl(n+\sqrt{n I(\nu\mid\gam)}\bigr)
\]
for an even isotropic log-concave probability measure \(\nu\) supported on
a symmetric convex body \(D\), with Poincar\'e constant at most \(P\) and
finite relative Fisher information.
A second-order Stein estimate and an information jump under convolution
give entropy decay faster than \(m^{-1}\) for normalised sums. Their
supports control the Gaussian measure of dilates, and a Brownian bridge
expansion gives the spectral estimate. A Gaussian-weighted ground state
supplies a measure whose Fisher information is bounded in terms of
\(\lambda_1(K)\). The square-root dependence then gives the bound of
order \(n\).
\end{abstract}

\maketitle


\section{Introduction}

Let \(n\ge1\) be an integer. A convex body \(K\subset\R^n\) is a
compact convex set with nonempty interior. Its volume and interior are
denoted by \(|K|\) and \(\Int K\), respectively. Let
\[
\lambda_1(K)
=\inf_{0\ne f\in H_0^1(\Int K)}
  \frac{\int_{\Int K}|\nabla f|^2\,dx}
       {\int_{\Int K}f^2\,dx}
\]
be the first eigenvalue of the Dirichlet Laplacian \(-\Delta\).
The same variational definition applies to a bounded open domain.
The affine reverse Faber--Krahn problem asks whether
\[
\inf_{S\in\SL_n(\R)}\lambda_1(SK)\,|K|^{2/n}\le Cn
\]
holds with an absolute constant. This question goes back to
Schmuckenschl\"ager~\cite{Sch11}; see also
\cite[Conjecture~5.1]{FP26}. An earlier bound of order \(n\log^2 n\)
is due to Klartag and E.~Milman; see
\cite[Proposition~4.8]{BK26} for the argument and its history.
Faifman and Polterovich~\cite[Theorem~5.2]{FP26} obtain the bound
\(Cn\log(2+n)\). We prove the conjectured bound of order \(n\).

\begin{theorem}\label{thm:affine-main}
There is a universal constant \(C>0\) such that every convex body
\(K\subset\R^n\) admits a map \(S\in\SL_n(\R)\) satisfying
\[
\lambda_1(SK)\,|K|^{2/n}\le Cn.
\]
\end{theorem}

For origin-symmetric bodies, the estimate holds in a classical affine
position defined by Gaussian measure.

Let \(\gam\) be standard Gaussian probability measure on \(\R^n\), with
density
\[
\varphi_n(x)=(2\pi)^{-n/2}e^{-|x|^2/2}.
\]
For a body \(K\), let
\[
\gamma_K(A)=\frac{\gam(A\cap K)}{\gam(K)}.
\]
An origin-symmetric body \(K\) of volume one is in
\emph{maximal Gaussian measure position} if
\[
\gam(K)=\max_{S\in\SL_n(\R)}\gam(SK).
\]
Such a position exists. Bobkov~\cite{Bob11} showed that the restricted
Gaussian has scalar covariance bounded above and below by absolute
constants in this position; see Section~\ref{sec:gaussian-position}.
Recent applications of this position to mean width and metric entropy
appear in Paouris and Pathak~\cite{PP26}.

\begin{theorem}\label{thm:gaussian-main}
There is a universal constant \(C>0\) such that every origin-symmetric
convex body \(K\subset\R^n\) with \(|K|=1\), in maximal Gaussian measure
position, satisfies
\[
\lambda_1(K)\le Cn.
\]
\end{theorem}

\subsection{Outline of the proof}

The Stein estimates used below originate in the spectral-gap
construction of Courtade, Fathi and Pananjady~\cite{CFP19} and its
higher-order extension by Fathi~\cite{Fathi21}. Their entropy estimates
build on the entropy--Stein--information method of Ledoux, Nourdin and
Peccati~\cite{LNP15}. We combine the second-order estimate with Bizeul's
information jump~\cite{Biz23}, a Fisher-information counterpart of the
entropy-jump inequality of Ball and Nguyen~\cite{BN12}.
Our central idea is to use the resulting entropy decay to bound
\(-\log\gam(tD)\) with a decay exponent greater than two, and then
estimate the survival probability of a Brownian bridge. This gives the
following spectral estimate.

\begin{theorem}\label{thm:information-spectrum}
Let \(D\subset\R^n\) be an origin-symmetric convex body, and let
\(\nu\) be an even isotropic log-concave probability measure supported
on \(D\). Suppose that
\[
\CP(\nu)\le P,\qquad P\ge2,
\]
and that \(I_0=I(\nu\mid\gam)<\infty\).
Then
\begin{equation}\label{eq:information-spectrum}
\lambda_1(D)\le C(P)\left(n+\sqrt{nI_0}\right),
\end{equation}
where \(C(P)\) is finite and depends only on \(P\).
\end{theorem}

For independent \(U_1,\ldots,U_m\) with law \(\nu\), let
\(\nu_m=\mathcal L(m^{-1/2}\sum_{j=1}^m U_j)\).
The second-order Stein construction in Section~\ref{sec:stein}, combined
with Bizeul's information jump, gives
\[
H(\nu_m\mid\gam)\le C(P)\sqrt{nI_0}\,m^{-1-\delta/2}
\]
for dyadic \(m\), where \(\delta=\delta(P)>0\).
Since \(\nu_m\) is supported on \(\sqrt mD\), this implies
\[
-\log\gam(tD)\le C(P)\sqrt{nI_0}\,t^{-2-\delta},\qquad t\ge1.
\]
The exponent greater than two allows the coefficients of a dyadic
Brownian bridge expansion to satisfy both a summable probability bound
and a summable uniform bound. The resulting survival estimate bounds the
Dirichlet heat kernel from below and proves
Theorem~\ref{thm:information-spectrum}.

To apply the theorem, we need a measure with bounded Poincar\'e constant
and controlled covariance, whose Fisher information may depend on the
eigenvalue. The restricted Gaussian in maximal Gaussian measure position
has the first two properties, but its density has a jump at the boundary.
In Section~\ref{sec:main-proof}, we convolve it with the probability
measure whose density is proportional to \(\psi^2e^{-|x|^2/2}\), where
\(\psi\) is a first Dirichlet eigenfunction. After covariance
normalisation, the resulting measure satisfies
\[
\CP(\nu)\le C,\qquad I(\nu\mid\gam)\le8\lambda_1(K)+n.
\]
Its support is a linear image of \(2K\), with uniformly bounded
distortion. Theorem~\ref{thm:information-spectrum} therefore gives
\[
\lambda_1(K)\le C\left(n+\sqrt{n\lambda_1(K)}\right).
\]
Absorbing the square-root term proves Theorem~\ref{thm:gaussian-main}.
Volume normalisation and the symmetric-intersection argument complete
the proof of Theorem~\ref{thm:affine-main}.

\subsection{Notation and conventions}

A probability measure is \emph{isotropic} if it has mean zero and
covariance \(\Id\). It is \emph{even} if it is invariant under
\(x\mapsto-x\). The law of a random vector \(U\) is denoted by
\(\mathcal L(U)\). Quantities such as \(\CP(U)\) and \(J(U)\) refer
to its law. A dyadic integer is \(m=2^r\) with \(r\in\{0,1,2,\ldots\}\).

For a probability density \(\rho\), define its ordinary Fisher
information by
\begin{equation}\label{eq:ordinary-fisher}
J(\rho\,dx)=
\begin{cases}
4\displaystyle\int_{\R^n}|\nabla\sqrt{\rho}|^2\,dx,
 &\sqrt{\rho}\in H^1(\R^n),\\
+\infty,&\text{otherwise}.
\end{cases}
\end{equation}
For smooth positive densities this equals
\(\int|\nabla\log\rho|^2\rho\,dx\).
For \(\nu=\rho\,dx\), its relative entropy with respect to \(\gam\) is
\[
H(\nu\mid\gam)=\int_{\R^n}
 \rho\log\frac{\rho}{\varphi_n}\,dx,
\]
with \(0\log0=0\); for a measure without a density we set
\(H(\nu\mid\gam)=+\infty\). For a density with finite second moment,
we use the
Sobolev definition
\[
I(\nu\mid\gam)=
\int_{\R^n}|2\nabla\sqrt\rho+x\sqrt\rho|^2\,dx
\]
when \(\sqrt\rho\in H^1(\R^n)\), and \(I(\nu\mid\gam)=+\infty\)
otherwise. For smooth positive densities this is
\(\int|\nabla\log(\rho/\varphi_n)|^2\rho\,dx\).
If \(J(\nu)<\infty\), integration by parts with cutoffs gives
\(\int x\cdot\nabla\rho\,dx=-n\), and hence
\[
I(\nu\mid\gam)=J(\nu)+\int|x|^2\,d\nu-2n.
\]
In particular, for isotropic \(\nu\),
\begin{equation}\label{eq:relative-ordinary}
I(\nu\mid\gam)=J(\nu)-n.
\end{equation}
The cutoff limit is justified by Cauchy--Schwarz, since
\(|x|\sqrt\rho\) and \(\nabla\sqrt\rho\) belong to \(L^2(dx)\).

The Poincar\'e constant \(\CP(\nu)\) is the least constant such that,
for every smooth function \(f\),
\begin{equation}\label{eq:poincare}
\Var_\nu f\le \CP(\nu)\int|\nabla f|^2\,d\nu.
\end{equation}
The inequality extends to Sobolev limits. In particular, isotropy implies
\(\CP(\nu)\ge1\).
For symmetric matrices, \(A\preceq B\) means that \(B-A\) is positive
semidefinite. We denote the Euclidean operator norm by
\(\|\cdot\|_{\op}\), and the Hilbert--Schmidt norm on matrices and
tensors by \(\|\cdot\|_{\HS}\). Thus
\(\|T\|_{L^2(\nu)}^2=\int\|T\|_{\HS}^2\,d\nu\).
The colon denotes coefficientwise contraction. For a vector field
\(F\), our derivative convention is \((DF)_{ij}=\partial_jF_i\).
For a third-order tensor \(T\),
\[
T:D^2F=\sum_{i,j,k}T_{ijk}\,\partial_j\partial_k F_i.
\]
For a matrix \(M\), the vector \(T:M\) has components
\((T:M)_i=\sum_{j,k}T_{ijk}M_{jk}\).

Symbols \(c,C>0\) denote absolute constants whose values may change from
one occurrence to the next. A constant denoted by \(C(P)\), or \(C(p)\),
depends only on the displayed parameter.

\section{Gaussian position and elementary analytic facts}
\label{sec:gaussian-position}

The maximal Gaussian measure position provides both a lower covariance
bound and a Poincar\'e inequality at a fixed scale.

\begin{proposition}[Bobkov's Gaussian position]\label{prop:gaussian-position}
Every origin-symmetric convex body \(K\subset\R^n\) of volume one has a
maximal Gaussian measure position obtained by a map in \(\SL_n(\R)\).
In that position,
\begin{equation}\label{eq:alpha-covariance}
\Cov(\gamma_K)=\alpha\Id,
\qquad c_0\le\alpha\le1,
\end{equation}
where \(c_0>0\) is an absolute constant.
\end{proposition}

This is a consequence of \cite[Proposition~3.1 and Lemma~4.1]{Bob11}.
We include the argument, since the covariance lower bound determines the
normalisation in the main proof.

\begin{proof}
The maximum exists. Indeed, rotational invariance reduces the
maximisation to positive definite maps with determinant one. If such
maps leave every compact subset of this matrix space, their smallest
singular values tend to zero along a subsequence. Since the original
body is bounded, the images lie in slabs of widths tending to zero.
Their Gaussian measures therefore tend to zero. Continuity of the
Gaussian measure under nonsingular linear maps now gives a maximiser.

At a maximiser, vary the body by \(e^{tA}\), where \(A\) is symmetric
and \(\tr A=0\). Differentiation gives
\[
0=-\int_K\langle Ax,x\rangle\,d\gam(x).
\]
Thus \(\int_K x\otimes x\,d\gam(x)\) is a scalar matrix. Symmetry gives
zero mean, so \(\Cov(\gamma_K)=\alpha\Id\).

A consequence of Milman's \(M\)-ellipsoid theorem \cite{Mil86} is that
there exists a volume-one ellipsoid \(E\) with
\[
|K\cap E|\ge c_1^n.
\]
The equivalent intersection formulation is also described in
\cite[Corollary~2.3]{Bob11}. Send \(E\) by a determinant-one map to the
Euclidean ball \(B\) of volume one. Its squared radius is at most \(Cn\);
hence in this position
\[
\gam(K)\ge \gam(K\cap B)
\ge (2\pi)^{-n/2}e^{-Cn}|K\cap B|
\ge e^{-C_1n}.
\]
The maximising position has at least this Gaussian measure.

Let \(q=\gam(K)\) in the maximising position and \(X\sim\gamma_K\).
Its differential entropy is
\[
h(X)=\frac n2\log(2\pi)+\frac{n\alpha}{2}+\log q.
\]
The Gaussian maximises entropy at fixed covariance, so
\[
h(X)\le\frac n2\log(2\pi e\alpha).
\]
Consequently
\[
-C_1n\le\log q
\le\frac n2(1+\log\alpha-\alpha),
\]
which implies \(\alpha\ge e^{-2C_1-1}\).
Finally, the restricted Gaussian is \(1\)-strongly log-concave.
The Brascamp--Lieb Poincar\'e inequality implies
\(\Cov(X)\preceq\Id\), giving \(\alpha\le1\).
\end{proof}

For \(\kappa>0\), a density is called \(\kappa\)-strongly log-concave if
it has the form
\[
\rho(x)=e^{-\kappa|x|^2/2-W(x)}
\]
up to normalisation, with \(W:\R^n\to(-\infty,+\infty]\) convex.
Convex support is allowed. The Brascamp--Lieb inequality
\cite{BL76} gives
\begin{equation}\label{eq:strong-lc-poincare}
\CP(\rho\,dx)\le\kappa^{-1},
\qquad
\Cov(\rho\,dx)\preceq\kappa^{-1}\Id.
\end{equation}
For extended-valued \(W\), approximate \(W\) by finite smooth convex
potentials and pass to the variance inequality for compactly supported
smooth functions. This proves the same bound on convex supports; cutoff
approximation then gives the covariance estimate. In the applications
below the potentials are even, so the means are zero.

The following facts describe the behaviour of Poincar\'e constants and
Fisher information under convolution and linear maps.

\begin{lemma}\label{lem:convolution-facts}
Let \(X,Y\) be independent random vectors.
\begin{enumerate}[label=\textup{(\roman*)},leftmargin=2.2em]
\item If their Poincar\'e constants are finite, then
\[
\CP(X+Y)\le\CP(X)+\CP(Y).
\]
For a linear map \(A\),
\[
\CP(AX)\le\norm{A}_{\op}^2\CP(X).
\]
Consequently, normalised sums of independent copies of a law with
Poincar\'e constant at most \(P\) also have Poincar\'e constant at most \(P\).

\item If \(Y\) has finite ordinary Fisher information, then
\[
J(X+Y)\le J(Y).
\]
Here \(X\) may have an arbitrary law.
For an invertible matrix \(A\),
\[
J(AX)\le\norm{A^{-1}}_{\op}^2 J(X).
\]

\item Let \(U\sim\nu\) and \(G\sim\gam\) be independent, with
\(\nu\) isotropic. For \(t\ge0\), define
\[
\nu_t=\mathcal L(e^{-t}U+\sqrt{1-e^{-2t}}\,G).
\]
Then
\[
\CP(\nu_t)\le e^{-2t}\CP(\nu)+1-e^{-2t}.
\]
If \(I(\nu\mid\gam)<\infty\), then
\[
I(\nu_t\mid\gam)\le e^{-2t}I(\nu\mid\gam).
\]
\end{enumerate}
\end{lemma}

\begin{proof}
For (i), apply the conditional variance decomposition and the
Poincar\'e inequalities to \(f(X+Y)\). The chain rule and Jensen's
inequality give the stated sum of constants. The linear-map statement
is another application of the chain rule.

For (ii), let \(g\) and \(h\) be the densities of \(Y\) and \(X+Y\).
Finite Fisher information
implies \(\sqrt g\in H^1\), and therefore
\(\nabla g=2\sqrt g\,\nabla\sqrt g\in L^1\).
Differentiating the convolution in the weak sense yields
\[
\nabla\log h(X+Y)
=\E[\nabla\log g(Y)\mid X+Y]
\]
on the set \(\{h>0\}\). Conditional Jensen's inequality proves the
information bound. Under \(x\mapsto Ax\), the score is transformed by
\(A^{-\mathsf T}\), which proves the second assertion.

The Poincar\'e assertion in (iii) follows from (i). The relative Fisher
contraction is the standard gradient commutation and
Cauchy--Schwarz estimate for the Ornstein--Uhlenbeck (OU) semigroup; see
\cite{BGL14}. It remains valid for finite Sobolev information by
approximation.
\end{proof}

\section{Second-order Stein kernels}
\label{sec:stein}

We use the spectral-gap construction of Stein kernels
\cite{CFP19} and its higher-order extension
\cite[Theorem~2.4]{Fathi21}. The argument below gives the second-order
identity for vector fields. Evenness makes the tensor centred, which
gives its decay under normalised summation.

Throughout this section, the density of \(\nu\) is smooth and strictly
positive on \(\R^n\), \(\nu\) is even and isotropic, and
\(\CP(\nu)\le P\). Let \(U\sim\nu\).
In the application, these assumptions will be obtained
by Ornstein--Uhlenbeck regularisation of a compactly supported measure.

Let \(H^1(\nu)\) be the closure of \(C_c^\infty(\R^n)\) in the norm
\[
\left(\int |f|^2\,d\nu+\int|\nabla f|^2\,d\nu\right)^{1/2}.
\]
The gradient is closable because the density is bounded above and
bounded away from zero on each compact set, so convergence in the
weighted spaces implies the local convergence needed to test weak
derivatives. Constants and coordinate functions belong to this space:
use smooth cutoffs \(\chi_R(x)=\chi(x/R)\) and \(x_j\chi_R(x)\).
The coordinate error tends to zero by the finite second moment, while
the extra derivative \(x_j\nabla\chi_R\) is uniformly bounded on an
annulus whose probability tends to zero. The Poincar\'e inequality makes
\(\int|\nabla f|^2\,d\nu\) a coercive quadratic form on its mean-zero
subspace.

\begin{lemma}\label{lem:second-stein}
There exists a third-order tensor \(T\in L^2(\nu)\), with
\[
\E T(U)=0,\qquad
\E\norm{T(U)}_{\HS}^2\le P(P-1)n,
\]
such that, for every \(F\in C_c^\infty(\R^n;\R^n)\),
\begin{equation}\label{eq:second-stein-identity}
\E[U\cdot F(U)-\Div F(U)]
=\E[T(U):D^2F(U)].
\end{equation}
The tensor \(T\) can be chosen odd.
\end{lemma}

\begin{proof}
On the mean-zero subspace of \(H^1(\nu)\), the functional
\(f\mapsto\int x_i f\,d\nu\) is continuous for the gradient norm, since
\[
\left|\int x_i f\,d\nu\right|
\le \norm{f}_{L^2(\nu)}
\le\sqrt P\,\norm{\nabla f}_{L^2(\nu)}
\]
when \(\int f\,d\nu=0\). The Lax--Milgram theorem gives a unique
mean-zero \(u_i\in H^1(\nu)\) such that
\begin{equation}\label{eq:first-poisson}
\int\nabla u_i\cdot\nabla f\,d\nu=\int x_i f\,d\nu
\qquad(f\in H^1(\nu)).
\end{equation}
Testing with \(u_i\) gives
\[
\int|\nabla u_i|^2\,d\nu\le P.
\]
Indeed, the left-hand side equals \(\int x_i u_i\,d\nu\), which is at
most \(\sqrt P\,\|\nabla u_i\|_{L^2(\nu)}\).
Define \(\tau_{ij}=\partial_j u_i\). Testing
\eqref{eq:first-poisson} with \(x_j\) gives
\(\E\tau_{ij}=\delta_{ij}\). Therefore
\begin{equation}\label{eq:first-stein-size}
\E\norm{\tau-\Id}_{\HS}^2
=\sum_i\int|\nabla u_i|^2\,d\nu-n
\le(P-1)n.
\end{equation}
The equation also gives the first-order Stein identity
\[
\E[U\cdot F(U)]=\E[\tau(U):DF(U)].
\]

The measure is even. Reflection of
\eqref{eq:first-poisson} and uniqueness show that \(u_i\) is odd.
Thus \(\tau\) is even. For each pair \(i,j\), solve the mean-zero weak
equation
\begin{equation}\label{eq:second-poisson}
\int\nabla v_{ij}\cdot\nabla f\,d\nu
=\int(\tau_{ij}-\delta_{ij})f\,d\nu.
\end{equation}
The source is mean zero and belongs to \(L^2(\nu)\), so
the Lax--Milgram theorem and the Poincar\'e inequality give
\[
\int|\nabla v_{ij}|^2\,d\nu
\le P\int|\tau_{ij}-\delta_{ij}|^2\,d\nu.
\]
The source is even; hence the mean-zero solution \(v_{ij}\) is even.
Set
\[
T_{ijk}=\partial_k v_{ij}.
\]
Then \(T\) is odd and centred, and summation together with
\eqref{eq:first-stein-size} yields
\[
\E\norm T_{\HS}^2\le P(P-1)n.
\]
Finally, apply \eqref{eq:second-poisson} with
\(f=\partial_j F_i\), sum over \(i,j\), and use the first-order
Stein identity. This gives \eqref{eq:second-stein-identity}.
\end{proof}

The following construction is the second-order case of
\cite[Lemma~4.1 and Corollary~4.2]{Fathi21}.

\begin{lemma}[Normalised sums]\label{lem:stein-sums}
Let \(U_1,\ldots,U_m\) be independent with law \(\nu\), and define
\[
W_m=m^{-1/2}\sum_{\ell=1}^m U_\ell.
\]
Then the tensor
\begin{equation}\label{eq:sum-tensor}
T_m(W_m)=m^{-3/2}
\E\left[\sum_{\ell=1}^m T(U_\ell)\,\middle|\,W_m\right]
\end{equation}
satisfies the second-order Stein identity for \(W_m\), and
\begin{equation}\label{eq:sum-tensor-bound}
\E\norm{T_m(W_m)}_{\HS}^2
\le\frac{P(P-1)n}{m^2}.
\end{equation}
\end{lemma}

\begin{proof}
Condition on all variables except \(U_\ell\), and apply
\eqref{eq:second-stein-identity} to the field obtained by composing \(F\)
with the normalised summation map. The first and second derivatives
contribute \(m^{-1/2}\) and \(m^{-1}\), respectively. Summing over
\(\ell\) gives the coefficient \(m^{-3/2}\) in
\eqref{eq:sum-tensor}.

Conditional Jensen's inequality and independence give
\[
\E\norm{T_m(W_m)}_{\HS}^2
\le m^{-3}\E\left\|\sum_{\ell=1}^m T(U_\ell)\right\|_{\HS}^2.
\]
Since \(T\) is centred, all cross terms vanish, and the last expression
is \(m^{-2}\E\norm T_{\HS}^2\). Apply
Lemma~\ref{lem:second-stein}.
\end{proof}

We also need the second-order entropy--Stein--information estimate. The
proof follows \cite[Theorem~3.1 and Lemma~3.2]{Fathi21}, in the form
suited to the vector-field identity above; see \cite{LNP15} for the
underlying semigroup method.

\begin{lemma}[Entropy from a second-order tensor]\label{lem:hsi}
Let \(U\) have an isotropic law \(\nu\) with finite relative Fisher
information and admit a tensor \(T\in L^2(\nu)\) satisfying
\eqref{eq:second-stein-identity}. Then
\begin{equation}\label{eq:hsi}
H(\nu\mid\gam)
\le\sqrt2\,\norm T_{L^2(\nu)}\sqrt{I(\nu\mid\gam)}.
\end{equation}
\end{lemma}

\begin{proof}
For \(t>0\), let \(a=e^{-t}\), \(b=\sqrt{1-e^{-2t}}\), and
\(U_t=aU+bG\), where
\(G\sim\gam\) is independent. Denote the relative score of its law by
\[
r_t=\nabla\log\frac{d\nu_t}{d\gam}.
\]
We first work against a vector field \(F\in C_c^\infty(\R^n;\R^n)\).
Integration by parts for the relative score gives
\[
\E[r_t(U_t)\cdot F(U_t)]
=\E[U_t\cdot F(U_t)-\Div F(U_t)].
\]
Gaussian integration by parts in the term containing \(bG\), followed
by \eqref{eq:second-stein-identity} in \(U\), gives
\begin{align*}
\E[U_t\cdot F(U_t)-\Div F(U_t)]
&=a\,\E[U\cdot F(aU+bG)]-a^2\E[\Div F(aU+bG)]\\
&=a^3\E[T(U):D^2F(aU+bG)].
\end{align*}
For every fixed value of \(G\), the field \(x\mapsto F(ax+bG)\)
is smooth and compactly supported, so the Stein identity applies.
Two further Gaussian integrations by parts give
\[
\E[r_t(U_t)\cdot F(U_t)]
=\frac{a^3}{b^2}
\E\bigl[(T(U):(G\otimes G-\Id))\cdot F(U_t)\bigr].
\]
The derivatives of \(F\) are bounded, and \(T\in L^2(\nu)\), which
justifies these integrations. Equivalently, the relative score is
identified weakly as
\begin{equation}\label{eq:ou-score}
r_t(U_t)
=\frac{a^3}{b^2}
\E[T(U):(G\otimes G-\Id)\mid U_t].
\end{equation}

For a fixed tensor \(T\), the Gaussian fourth-moment formula gives
\begin{equation}\label{eq:hermite-contraction}
\E_G|T:(G\otimes G-\Id)|^2
=\sum_{i,j,k}\bigl(T_{ijk}^2+T_{ijk}T_{ikj}\bigr)
\le2\norm T_{\HS}^2.
\end{equation}
Thus \eqref{eq:ou-score} and conditional Jensen's inequality imply
\begin{equation}\label{eq:ou-second-bound}
I(\nu_t\mid\gam)
\le\frac{2e^{-6t}}{(1-e^{-2t})^2}\norm T_{L^2(\nu)}^2.
\end{equation}
Together with the relative Fisher-information contraction under the OU
semigroup, this gives
\[
I(\nu_t\mid\gam)
\le
\min\left\{
e^{-2t}I(\nu\mid\gam),
\frac{2e^{-6t}}{(1-e^{-2t})^2}\norm T_{L^2(\nu)}^2
\right\}.
\]

The entropy dissipation identity for the OU semigroup is
\[
H(\nu\mid\gam)=\int_0^\infty I(\nu_t\mid\gam)\,dt;
\]
see, for example, \cite{BGL14}. Finite information ensures finite
entropy by the Gaussian logarithmic Sobolev inequality. The same
inequality gives
\(H(\nu_t\mid\gam)\le I(\nu_t\mid\gam)/2\to0\) as \(t\to\infty\).
At \(t\downarrow0\), lower semicontinuity of entropy and its contraction
under the semigroup give \(H(\nu_t\mid\gam)\to H(\nu\mid\gam)\).
Thus both endpoints in the dissipation identity are justified. Set
\(I_*=I(\nu\mid\gam)\) and \(s=\norm T_{L^2(\nu)}\).
If \(I_*=0\), the Gaussian logarithmic Sobolev inequality gives
\(H(\nu\mid\gam)=0\).
If \(s=0\), \eqref{eq:ou-second-bound} and entropy dissipation give
\(H(\nu\mid\gam)=0\). We may therefore assume \(I_*,s>0\). With
\(z=1-e^{-2t}\),
\[
H(\nu\mid\gam)
\le\frac12\int_0^1
 \min\left\{I_*,\frac{2s^2(1-z)^2}{z^2}\right\}\,dz
\le\frac12\int_0^1
 \min\left\{I_*,\frac{2s^2}{z^2}\right\}\,dz.
\]
If \(\sqrt2s/\sqrt{I_*}\le1\), split the integral at this value. The
result is at most \(\sqrt2s\sqrt{I_*}\). If
\(\sqrt2s/\sqrt{I_*}>1\), the integral is at most \(I_*/2\), which
satisfies the same upper bound. This proves
\eqref{eq:hsi}.
\end{proof}

\section{Entropy decay and Gaussian measure}
\label{sec:entropy}

The Stein estimate supplies one power of \(m^{-1}\) in the entropy of a
normalised sum. Bizeul's information jump supplies the additional decay
needed for the Brownian bridge argument.

\begin{proposition}[Bizeul's information jump]\label{prop:fisher-jump}
Let \(U,U'\) be independent and identically distributed, with an
isotropic log-concave law \(\nu\) and
\(\CP(\nu)\le P\). If \(I(\nu\mid\gam)<\infty\), then
\begin{equation}\label{eq:fisher-jump}
I\left(\mathcal L\left(\frac{U+U'}{\sqrt2}\right)\middle|\gam\right)
\le \left(1-\frac1{8P}\right)I(\nu\mid\gam).
\end{equation}
\end{proposition}

\begin{proof}
This is \cite[Theorem~2]{Biz23}, with equal input laws and mixing
parameter \(1/2\). The deficit in that theorem is expressed in ordinary
Fisher information, while its lower bound uses relative Fisher
information. The output and input laws are isotropic, so
\eqref{eq:relative-ordinary} shows that their ordinary and relative
information deficits coincide. The resulting lower bound on the
deficit is \(I(\nu\mid\gam)/(8P)\).

For a nonsmooth density, apply the estimate first to
\(\nu_\eps=\mathcal L(e^{-\eps}U+\sqrt{1-e^{-2\eps}}\,G)\).
This law is smooth, positive, isotropic and log-concave, and
Lemma~\ref{lem:convolution-facts} gives
\(\CP(\nu_\eps)\le P\) and
\(I(\nu_\eps\mid\gam)\le I(\nu\mid\gam)\).
Taking a normalised sum of two independent copies commutes with this
regularisation. Lower semicontinuity of relative Fisher information at
the output as \(\eps\downarrow0\) proves \eqref{eq:fisher-jump}.
\end{proof}

Let \(I_0=I(\nu\mid\gam)<\infty\), and let \(\nu_m\) be the law of
the normalised sum of \(m\) independent copies with law \(\nu\).
Log-concavity follows
from the Pr\'ekopa--Leindler theorem, isotropy is preserved, and
Lemma~\ref{lem:convolution-facts} gives \(\CP(\nu_m)\le P\).
Consequently, for dyadic \(m\),
\begin{equation}\label{eq:dyadic-fisher}
I(\nu_m\mid\gam)\le m^{-\delta}I_0,
\qquad
\delta=-\log_2\left(1-\frac1{8P}\right)>0.
\end{equation}

\begin{proposition}[Entropy of symmetric normalised sums]
\label{prop:entropy-clt}
Suppose \(\nu\) is an even isotropic log-concave probability supported
on a convex body, with \(\CP(\nu)\le P\), \(P\ge2\), and finite
\(I_0=I(\nu\mid\gam)\). Then, for dyadic \(m\),
\begin{equation}\label{eq:entropy-clt}
H(\nu_m\mid\gam)
\le \sqrt{2P(P-1)n I_0}\,m^{-1-\delta/2},
\end{equation}
where \(\delta\) is defined in \eqref{eq:dyadic-fisher}.
\end{proposition}

\begin{proof}
Assume first that \(\nu\) has a smooth positive density on \(\R^n\),
is even, isotropic and log-concave, and satisfies the stated Poincar\'e
and information bounds.
Lemmas~\ref{lem:stein-sums} and \ref{lem:hsi}, together with
\eqref{eq:dyadic-fisher}, give
\[
H(\nu_m\mid\gam)
\le \sqrt2\,\norm{T_m}_{L^2(\nu_m)}
       \sqrt{I(\nu_m\mid\gam)}
\le \sqrt{2P(P-1)n I_0}\,m^{-1-\delta/2}.
\]

For the compactly supported law in the statement, set
\[
\nu_\eps
=\mathcal L(e^{-\eps}U+\sqrt{1-e^{-2\eps}}\,G).
\]
The density of \(\nu_\eps\) is smooth, strictly positive, even, and
log-concave. It remains isotropic, and
Lemma~\ref{lem:convolution-facts} gives
\[
\CP(\nu_\eps)\le e^{-2\eps}P+1-e^{-2\eps}\le P,
\qquad
I(\nu_\eps\mid\gam)\le I_0.
\]
Since the original law has compact support, its Gaussian convolution
has Gaussian tails and all moments are finite. Thus the Sobolev
cutoffs and integrations used in Section~\ref{sec:stein} apply.

For every fixed dyadic \(m\), normalised summation commutes with OU
regularisation:
\[
(\nu_\eps)_m=(\nu_m)_\eps.
\]
The preceding entropy estimate is uniform in \(\eps\).
As \(\eps\downarrow0\), weak convergence and lower semicontinuity of
relative entropy therefore yield
\[
H(\nu_m\mid\gam)
\le\liminf_{\eps\downarrow0}H((\nu_\eps)_m\mid\gam)
\le\sqrt{2P(P-1)n I_0}\,m^{-1-\delta/2}.
\]
\end{proof}

\begin{lemma}[Entropy and support]\label{lem:support-entropy}
If a probability measure \(\eta\) is supported on a measurable set \(A\)
of positive Gaussian measure, then
\[
-\log\gam(A)\le H(\eta\mid\gam).
\]
\end{lemma}

\begin{proof}
If the relative entropy is infinite there is nothing to prove.
Otherwise \(\eta\ll\gam(\cdot\mid A)\), and
\[
H(\eta\mid\gam)
=H(\eta\mid\gam(\cdot\mid A))-\log\gam(A)
\ge-\log\gam(A).
\]
\end{proof}

\begin{proposition}\label{prop:gaussian-profile}
Under the assumptions of Theorem~\ref{thm:information-spectrum}, let
\(p=2+\delta>2\). There is
\[
A\le C(P)\sqrt{n I_0}
\]
such that
\begin{equation}\label{eq:gaussian-profile}
-\log\gam(tD)\le A t^{-p}\qquad(t\ge1).
\end{equation}
\end{proposition}

\begin{proof}
Convexity gives
\(\supp(\nu_m)\subset m^{-1/2}(D+\cdots+D)=\sqrt mD\).
Apply Lemma~\ref{lem:support-entropy} and
Proposition~\ref{prop:entropy-clt} to obtain, for dyadic \(m\),
\[
-\log\gam(\sqrt mD)
\le \sqrt{2P(P-1)n I_0}\,m^{-p/2}.
\]
For a given \(t\ge1\), choose the largest dyadic integer \(m\le t^2\).
Then \(t^2/2<m\le t^2\). Monotonicity under dilation gives
\[
-\log\gam(tD)
\le 2^{p/2}\sqrt{2P(P-1)n I_0}\,t^{-p}.
\]
This proves the assertion.
\end{proof}

\section{Brownian bridges and the first eigenvalue}
\label{sec:bridge}

The preceding section gives a bound on Gaussian measure under dilation.
We now convert that bound into a lower bound for Brownian bridge
survival, and hence an upper bound for the first Dirichlet eigenvalue.

\begin{lemma}
\label{lem:bridge}
Let \(D\subset\R^n\) be an origin-symmetric convex body. Suppose that
\(p>2\) and \(A\ge0\) satisfy
\[
-\log\gam(tD)\le A t^{-p}\qquad(t\ge1).
\]
Then
\begin{equation}\label{eq:bridge-eigenvalue}
\lambda_1(D)\le C(p)(A+n).
\end{equation}
\end{lemma}

\begin{proof}
Let \(\beta\) be the standard \(n\)-dimensional Brownian bridge from
zero to zero on \([0,1]\), associated with Brownian motion of generator
\(\Delta/2\). Its L\'evy--Ciesielski expansion is
\begin{equation}\label{eq:bridge-expansion}
\beta(t)=\sum_{j=0}^\infty 2^{-j/2-1}
 \sum_{k=0}^{2^j-1}G_{jk}h_{jk}(t),
\end{equation}
where the \(G_{jk}\) are independent standard Gaussian vectors.
The function \(h_{jk}\) is supported on
\([k2^{-j},(k+1)2^{-j}]\), is linear on each half of that interval,
vanishes at its endpoints, and equals one at the midpoint.
At each level \(j\) and time \(t\), at most one tent is nonzero.

The coefficients are the normalised midpoint residuals
\[
2^{j/2+1}\left[
 \beta\bigl((2k+1)2^{-j-1}\bigr)
 -\frac12\left\{\beta(k2^{-j})+\beta((k+1)2^{-j})\right\}
\right].
\]
Conditional midpoint variances are \(2^{-j-2}\), and Gaussian
conditioning gives independence of the residuals. This yields
\eqref{eq:bridge-expansion}; the standard L\'evy construction is discussed
in \cite[Section~1.1.2]{MP10}.

Choose a fixed number \(a\) satisfying
\begin{equation}\label{eq:choice-a}
\frac1p<a<\frac12.
\end{equation}
Set \(b_j=2^{aj}\) and consider the coefficient event
\[
E=\bigcap_{j\ge0}\ \bigcap_{k=0}^{2^j-1}\{G_{jk}\in b_jD\}.
\]
By independence,
\begin{align}
-\log\Prob(E)
&=\sum_{j=0}^\infty 2^j[-\log\gam(b_jD)]\notag\\
&\le A\sum_{j=0}^\infty2^{(1-ap)j}
=\frac{A}{1-2^{1-ap}}.
\label{eq:bridge-probability}
\end{align}
For finite intersections, independence gives the corresponding finite
product. Continuity of probability under decreasing intersections gives
the displayed identity for \(E\); the product is positive because
\(ap>1\).

Let \(\norm{\cdot}_D\) be the norm with unit ball \(D\).
On \(E\), \eqref{eq:bridge-expansion} gives the deterministic uniform
bound
\begin{equation}\label{eq:bridge-radius}
\sup_{t\in[0,1]}\norm{\beta(t)}_D
\le\frac12\sum_{j=0}^\infty2^{(a-1/2)j}
=R,\qquad
R=\frac1{2(1-2^{a-1/2})}.
\end{equation}
This same convergent series controls uniform convergence of the
expansion on \(E\). The expansion agrees almost surely with the
continuous bridge.

Let \(\Omega=2R \Int D\).
Since \(0\in \Int D\), the compact set \(RD\) is contained in \(\Omega\)
with a positive distance from its boundary. Thus \(E\) implies
survival of the bridge inside the open domain. If
\[
q=\Prob\{\beta(t)\in\Omega\text{ for every }t\in[0,1]\},
\]
then
\begin{equation}\label{eq:q-lower}
q\ge\exp\left(-\frac{A}{1-2^{1-ap}}\right).
\end{equation}

Let \(p_t^\Omega\) be the Dirichlet heat kernel for the generator
\(\Delta/2\). The killed Brownian transition representation and bridge
conditioning give
\begin{equation}\label{eq:bridge-kernel}
p_1^\Omega(0,0)=q\,p_1^{\R^n}(0,0).
\end{equation}
See \cite[Section~7.4]{MP10} for the killed-process representation.

Let \(\lambda_\ell(\Omega)\) and \(\phi_\ell\) be Dirichlet
eigenvalues of \(-\Delta\) and an orthonormal eigenbasis. Positivity of
the squared coefficients in the diagonal spectral expansion yields
\begin{align}
p_1^\Omega(0,0)
&=\sum_\ell e^{-\lambda_\ell(\Omega)/2}\phi_\ell(0)^2\notag\\
&\le e^{-\lambda_1(\Omega)/4}
  \sum_\ell e^{-\lambda_\ell(\Omega)/4}\phi_\ell(0)^2\notag\\
&=e^{-\lambda_1(\Omega)/4}p_{1/2}^\Omega(0,0)\notag\\
&\le e^{-\lambda_1(\Omega)/4}p_{1/2}^{\R^n}(0,0).
\label{eq:heat-kernel-bound}
\end{align}
The heat kernel and eigenfunctions are smooth in the interior of
\(\Omega\), so these diagonal evaluations are well defined. The last
inequality is domination by the free heat kernel.
For Brownian motion of generator \(\Delta/2\),
\(p_t^{\R^n}(0,0)=(2\pi t)^{-n/2}\), so
\eqref{eq:bridge-kernel}--\eqref{eq:heat-kernel-bound} give
\[
q\le 2^{n/2}e^{-\lambda_1(\Omega)/4}.
\]
It follows that
\[
\lambda_1(\Omega)
\le4\log(1/q)+2n\log2
\le\frac{4A}{1-2^{1-ap}}+2n\log2.
\]
Finally,
\[
\lambda_1(D)=(2R)^2\lambda_1(\Omega).
\]
Choosing, for example,
\(a=\tfrac12(\tfrac1p+\tfrac12)\) makes every resulting constant
depend only on \(p\), which proves \eqref{eq:bridge-eigenvalue}.
\end{proof}

\begin{proof}[Proof of Theorem~\ref{thm:information-spectrum}]
Proposition~\ref{prop:gaussian-profile} gives
\[
-\log\gam(tD)\le C(P)\sqrt{nI_0}\,t^{-p},
\qquad
p=2-\log_2\left(1-\frac1{8P}\right)>2.
\]
Apply Lemma~\ref{lem:bridge}. Both the exponent \(p\) and the
constants depend only on \(P\), yielding
\eqref{eq:information-spectrum}.
\end{proof}

\begin{remark}\label{rem:critical-exponent}
The two series in \eqref{eq:bridge-probability} and
\eqref{eq:bridge-radius} converge simultaneously when
\(1/p<a<1/2\), which is possible exactly when \(p>2\).
The centred Stein tensor contributes a factor \(m^{-1}\) to the entropy
estimate. The information jump adds \(m^{-\delta/2}\), giving the
Gaussian dilation exponent \(2+\delta\).
\end{remark}

\section{Proof of the affine reverse Faber--Krahn inequality}
\label{sec:main-proof}

Theorem~\ref{thm:information-spectrum} reduces the spectral estimate to
the construction of a measure with controlled Poincar\'e constant and
Fisher information. The ground state provides the required boundary
decay; convolution with the restricted Gaussian supplies a covariance
lower bound.

Let \(K\) be origin-symmetric, of volume one, in maximal Gaussian
measure position. Let
\[
\lambda=\lambda_1(K),\qquad
\Cov(\gamma_K)=\alpha\Id,\qquad c_0\le\alpha\le1.
\]
Let \(\psi\in H_0^1(\Int K)\) be a positive first eigenfunction.
It is unique up to a positive scalar. Since \(K=-K\), it can be
chosen even. The Brascamp--Lieb log-concavity theorem for the first
Dirichlet eigenfunction on a convex domain \cite{BL76} gives
log-concavity of \(\psi\). The extension by zero is understood in this
log-concavity statement.

\begin{lemma}[Gaussian-weighted ground state]\label{lem:groundstate}
Define
\[
v(x)=\psi(x)e^{-|x|^2/4}\quad(x\in \Int K),
\]
and extend \(v\) by zero. Let \(Y\) have probability density
\[
q_Y(x)=\frac{v(x)^2}{\int_{\R^n}v^2}.
\]
Then \(Y\) is even and \(1\)-strongly log-concave, and
\[
\CP(Y)\le1,\qquad
\Cov(Y)\preceq\Id,\qquad
J(Y)\le4\lambda+n.
\]
More precisely,
\begin{equation}\label{eq:groundstate-exact}
J(Y)=4\lambda+\E|Y|^2.
\end{equation}
\end{lemma}

\begin{proof}
The density is proportional to
\(\psi^2e^{-|x|^2/2}\), so log-concavity of \(\psi\) gives
\(1\)-strong log-concavity. Symmetry gives zero mean.
Equation~\eqref{eq:strong-lc-poincare} supplies the Poincar\'e and
covariance bounds.

To compute the information, test the weak equation
\(-\Delta\psi=\lambda\psi\) with
\(\psi e^{-|x|^2/2}\). This is an admissible member of
\(H_0^1(\Int K)\). With \(w=e^{-|x|^2/2}\), this gives
\[
\int_K w|\nabla\psi|^2
-\int_K w\psi\,x\cdot\nabla\psi
=\lambda\int_K w\psi^2.
\]
On the other hand,
\[
|\nabla v|^2
=w|\nabla\psi|^2-w\psi\,x\cdot\nabla\psi
+\frac14|x|^2w\psi^2.
\]
Thus
\begin{equation}\label{eq:groundstate-energy}
\int|\nabla v|^2
=\lambda\int v^2+\frac14\int|x|^2v^2.
\end{equation}
Multiplication by the smooth Gaussian weight preserves \(H_0^1\), and
zero extension places \(v\) in \(H^1(\R^n)\). Therefore the definition
\eqref{eq:ordinary-fisher} applies and gives
\[
J(Y)=4\frac{\int|\nabla v|^2}{\int v^2}
=4\lambda+\E|Y|^2.
\]
The covariance bound and evenness imply \(\E|Y|^2\le n\).
\end{proof}

Independently of \(Y\), let \(X\sim\gamma_K\), and define
\begin{equation}\label{eq:constructed-measure}
\Sigma=\Cov(X+Y)=\alpha\Id+\Cov(Y),\qquad
Z=\Sigma^{-1/2}(X+Y),\qquad
D=2\Sigma^{-1/2}K.
\end{equation}
By Proposition~\ref{prop:gaussian-position} and
Lemma~\ref{lem:groundstate},
\begin{equation}\label{eq:sigma-bounds}
\alpha\Id\preceq\Sigma\preceq2\Id.
\end{equation}

\begin{lemma}\label{lem:constructed-bounds}
The law \(\nu=\mathcal L(Z)\) is even, isotropic, log-concave and
supported on \(D\). It satisfies
\begin{equation}\label{eq:constructed-poincare}
\CP(\nu)\le\frac2\alpha\le\frac2{c_0},
\end{equation}
and
\begin{equation}\label{eq:constructed-fisher}
I(\nu\mid\gam)\le8\lambda+n.
\end{equation}
\end{lemma}

\begin{proof}
Evenness, log-concavity and support follow from independence,
convolution of log-concave densities, and \(K+K=2K\).
The definition of \(\Sigma\) gives isotropy.
The law of \(X\) is \(1\)-strongly log-concave, so
\(\CP(X)\le1\). Lemma~\ref{lem:groundstate} gives
\(\CP(Y)\le1\). Hence Lemma~\ref{lem:convolution-facts} implies
\[
\CP(Z)\le\norm{\Sigma^{-1/2}}_{\op}^2\CP(X+Y)
\le\frac2\alpha.
\]

The same lemma, now for Fisher information, gives
\[
J(Z)\le\norm{\Sigma}_{\op}J(X+Y)
\le2J(Y)\le8\lambda+2n.
\]
Here only \(Y\) is required to have finite Fisher information; the
density of \(X\) may have a jump at \(\partial K\).
Since \(Z\) is isotropic, \eqref{eq:relative-ordinary} proves
\eqref{eq:constructed-fisher}.
\end{proof}

\begin{lemma}\label{lem:linear-rayleigh}
For any invertible \(A:\R^n\to\R^n\) and any bounded domain \(U\),
\[
\lambda_1(U)\le\norm A_{\op}^2\lambda_1(AU).
\]
\end{lemma}

\begin{proof}
For \(f\in H_0^1(AU)\), use \(f\circ A\in H_0^1(U)\) as a test
function. Change of variables gives
\[
\frac{\int_U|\nabla(f\circ A)|^2}{\int_U|f\circ A|^2}
=
\frac{\int_{AU}|A^{\mathsf T}\nabla f|^2}{\int_{AU}|f|^2}
\le\norm A_{\op}^2
\frac{\int_{AU}|\nabla f|^2}{\int_{AU}|f|^2}.
\]
Take the infimum over nonzero \(f\).
\end{proof}

\begin{proof}[Proof of Theorem~\ref{thm:gaussian-main}]
Use the measure in \eqref{eq:constructed-measure}.
The fixed number \(P_0=2/c_0\) is at least two and is universal.
Theorem~\ref{thm:information-spectrum} and
Lemma~\ref{lem:constructed-bounds} give
\[
\lambda_1(D)
\le C(P_0)\left(n+\sqrt{n(8\lambda+n)}\right).
\]
For \(A=2\Sigma^{-1/2}\), equation~\eqref{eq:sigma-bounds} gives
\(\norm A_{\op}^2\le4/\alpha\le4/c_0\).
Since \(D=AK\), Lemma~\ref{lem:linear-rayleigh} yields
\[
\lambda
\le C\left(n+\sqrt{n(8\lambda+n)}\right)
\le C'\left(n+\sqrt{n\lambda}\right),
\]
with absolute constants. By Young's inequality,
\[
C'\sqrt{n\lambda}
\le\frac{\lambda}{2}+\frac{(C')^2n}{2}.
\]
Absorbing \(\lambda/2\) proves
\[
\lambda_1(K)\le\bigl(2C'+(C')^2\bigr)n.
\]
\end{proof}

\subsection{Reduction to symmetric bodies}
\label{sec:general-case}

We use the classical bound that every convex body \(K\) contains a
centrally symmetric convex body of volume at least \(2^{-n}|K|\);
see \cite[Section~1.2]{HSTV22}. The short averaging argument gives the
required reduction.

\begin{proof}[Proof of Theorem~\ref{thm:affine-main}]
Suppose first that \(K=-K\). Let \(V=|K|\) and \(K_0=V^{-1/n}K\).
Choose \(S\in\SL_n(\R)\) placing \(K_0\) in maximal Gaussian measure
position. Theorem~\ref{thm:gaussian-main} gives
\[
Cn\ge\lambda_1(SK_0)
=V^{2/n}\lambda_1(SK),
\]
which proves the theorem for origin-symmetric bodies.

Now let \(K\) be an arbitrary convex body and \(V=|K|\).
By Fubini,
\[
\int_{\R^n}|K\cap(x-K)|\,dx=V^2.
\]
The integrand vanishes outside \(K+K=2K\), whose volume is \(2^nV\).
Therefore some \(x\in\R^n\) satisfies
\[
|K\cap(x-K)|\ge2^{-n}V.
\]
Set \(L=K\cap(x-K)\). It is a full-dimensional convex body,
centrally symmetric about \(x/2\), and \(L\subset K\).
The translated body \(L_0=L-x/2\) is origin-symmetric.

Apply the symmetric case to \(L_0\), obtaining
\(S\in\SL_n(\R)\) with
\(\lambda_1(SL_0)|L|^{2/n}\le Cn\).
Translation invariance gives
\(\lambda_1(SL_0)=\lambda_1(SL)\), and domain monotonicity gives
\(\lambda_1(SK)\le\lambda_1(SL)\). Hence
\[
\lambda_1(SK)V^{2/n}
\le Cn\left(\frac{V}{|L|}\right)^{2/n}
\le4Cn.
\]
Absorb the factor four into the universal constant.
\end{proof}

\begin{thebibliography}{99}

\bibitem{BGL14}
D.~Bakry, I.~Gentil, and M.~Ledoux,
\emph{Analysis and Geometry of Markov Diffusion Operators},
Grundlehren der mathematischen Wissenschaften, vol.~348,
Springer, Cham, 2014.
\href{https://doi.org/10.1007/978-3-319-00227-9}
{doi:10.1007/978-3-319-00227-9}.

\bibitem{BN12}
K.~Ball and V.~H.~Nguyen,
Entropy jumps for isotropic log-concave random vectors and spectral gap,
\emph{Studia Mathematica} \textbf{213} (2012), no.~1, 81--96.
\href{https://doi.org/10.4064/sm213-1-6}
{doi:10.4064/sm213-1-6}.

\bibitem{Biz23}
P.~Bizeul,
Entropy and information jump for log-concave vectors,
\emph{Comptes Rendus. Math\'ematique} \textbf{361} (2023), no.~G2,
487--493.
\href{https://doi.org/10.5802/crmath.390}
{doi:10.5802/crmath.390}.

\bibitem{BK26}
P.~Bizeul and B.~Klartag,
\emph{Distances between non-symmetric convex bodies: optimal bounds up to polylog},
arXiv:2510.20511v3 [math.MG], 2026.
\href{https://arxiv.org/abs/2510.20511v3}{arXiv:2510.20511v3}.

\bibitem{Bob11}
S.~G.~Bobkov,
On Milman's ellipsoids and \(M\)-position of convex bodies,
in \emph{Concentration, Functional Inequalities and Isoperimetry},
C.~Houdr\'e, M.~Ledoux, E.~Milman, and M.~Milman, eds.,
Contemporary Mathematics, vol.~545,
American Mathematical Society, Providence, RI, 2011, pp.~23--33.
\href{https://doi.org/10.1090/conm/545/10762}
{doi:10.1090/conm/545/10762}.

\bibitem{BL76}
H.~J.~Brascamp and E.~H.~Lieb,
On extensions of the Brunn--Minkowski and Pr\'ekopa--Leindler theorems,
including inequalities for log concave functions, and with an
application to the diffusion equation,
\emph{Journal of Functional Analysis} \textbf{22} (1976), no.~4,
366--389.
\href{https://doi.org/10.1016/0022-1236(76)90004-5}
{doi:10.1016/0022-1236(76)90004-5}.

\bibitem{CFP19}
T.~A.~Courtade, M.~Fathi, and A.~Pananjady,
Existence of Stein kernels under a spectral gap, and discrepancy bounds,
\emph{Annales de l'Institut Henri Poincar\'e, Probabilit\'es et Statistiques}
\textbf{55} (2019), no.~2, 777--790.
\href{https://doi.org/10.1214/18-AIHP898}{doi:10.1214/18-AIHP898}.

\bibitem{FP26}
D.~Faifman and I.~Polterovich,
\emph{The Faber--Krahn position of convex bodies and Gaussian measure
inequalities},
arXiv:2607.21539v2 [math.SP], revised 5 August 2026.
\href{https://arxiv.org/abs/2607.21539v2}
{arXiv:2607.21539v2}.

\bibitem{Fathi21}
M.~Fathi,
Higher-order Stein kernels for Gaussian approximation,
\emph{Studia Mathematica} \textbf{256} (2021), no.~3,
241--258.
\href{https://doi.org/10.4064/sm190415-28-10}
{doi:10.4064/sm190415-28-10}.

\bibitem{HSTV22}
H.~Huang, B.~A.~Slomka, T.~Tkocz, and B.-H.~Vritsiou,
Improved bounds for Hadwiger's covering problem via thin-shell estimates,
\emph{Journal of the European Mathematical Society}
\textbf{24} (2022), no.~4, 1431--1448.
\href{https://doi.org/10.4171/JEMS/1132}{doi:10.4171/JEMS/1132}.

\bibitem{LNP15}
M.~Ledoux, I.~Nourdin, and G.~Peccati,
Stein's method, logarithmic Sobolev and transport inequalities,
\emph{Geometric and Functional Analysis} \textbf{25} (2015), no.~1, 256--306.
\href{https://doi.org/10.1007/s00039-015-0312-0}
{doi:10.1007/s00039-015-0312-0}.

\bibitem{Mil86}
V.~D.~Milman,
An inverse form of the Brunn--Minkowski inequality, with applications
to the local theory of normed spaces,
\emph{Comptes Rendus de l'Acad\'emie des Sciences de Paris,
S\'erie I Math\'ematique} \textbf{302} (1986), no.~1,
25--28.

\bibitem{MP10}
P.~M\"orters and Y.~Peres,
\emph{Brownian Motion},
Cambridge Series in Statistical and Probabilistic Mathematics,
vol.~30, Cambridge University Press, Cambridge, 2010.
\href{https://doi.org/10.1017/CBO9780511750489}
{doi:10.1017/CBO9780511750489}.

\bibitem{PP26}
G.~Paouris and R.~Pathak,
\emph{Optimal mean width and metric entropy estimates for convex bodies},
arXiv:2607.29522v2 [math.MG], 2026.
\href{https://arxiv.org/abs/2607.29522v2}{arXiv:2607.29522v2}.

\bibitem{Sch11}
M.~Schmuckenschl\"ager,
Inequalities for exit times and eigenvalues of balls,
\emph{Potential Analysis} \textbf{35} (2011), no.~3, 287--300.
\href{https://doi.org/10.1007/s11118-010-9213-x}
{doi:10.1007/s11118-010-9213-x}.

\end{thebibliography}

\end{document}

